% GENERATED FILE. Edit canonical lessons, then run npm run generate:books.
% canonical-source-hash: fnv1a64:44c908ccefe578cc
% canonical-lessons: TA-C108-eaten, TA-C108-longtime, TA-C108-whatsnew, TA-C108-areyouwell, TA-C108-howareyou, TA-R108-second-pass-leave-greet, TA-R108-second-pass-care

\chapter{Greetings People Actually Say}
\label{ch:ta-greetings-people-say}
\hlchaptermodality{108}

\begin{chapteropening}
\textbf{By the end of this chapter:} \emph{I can greet someone the way Tamil speakers do day to day: have you eaten, long time, what's new, are you well, how are you.}

Complete the last lesson of chapter 108: i can greet someone the way Tamil speakers do day to day: have you eaten, long time, what's new, are you well, how are you.
\end{chapteropening}

\section[cāppiṭṭīrkaḷā?]{\ta{சாப்பிட்டீர்களா}? (cāppiṭṭīrkaḷā?) — have you eaten?}

\label{lesson:TA-C108-eaten}



\begin{quote}
Before the new one: say the Tamil for a good journey, then the Tamil for come again.
\end{quote}

\subsection*{You'll want to know: \ta{சாப்பிட்டீர்களா}?}
\textbf{\ta{சாப்பிட்டீர்களா}?} (\emph{cāppiṭṭīrkaḷā?}) — \textquotedblleft{}have you eaten?\textquotedblright{}.

Around a mealtime it is a greeting and a kindness, built on \textbf{\ta{சாப்பிடு}}, \textquotedblleft{}eat\textquotedblright{}.

\begin{grammarlens}[title={What you've built}]
A greeting about food.
\end{grammarlens}

\subsection*{Your turn}
\begin{itemize}
\raggedright
  \item \emph{Say it:} \emph{cāppiṭṭīrkaḷā?}
  \item \emph{Say it:} \emph{cāppiṭṭīrkaḷā?}, once more
  \item \emph{Say it:} ask \emph{cāppiṭṭīrkaḷā?}
  \item \emph{Recall:} say the Tamil for a good journey, then the Tamil for come again, then say \emph{cāppiṭṭīrkaḷā?} again
\end{itemize}

\begin{tcolorbox}[breakable,colback=teal!4,colframe=teal!35!black,title={Before you move on}]
What does \ta{சாப்பிட்டீர்களா}? mean? (\textquotedblleft{}Have you eaten?\textquotedblright{}.) Say it once more.
\end{tcolorbox}
\section[rompa nāḷāccu]{\ta{ரொம்ப} \ta{நாளாச்சு} (rompa nāḷāccu) — it's been a long time}

\label{lesson:TA-C108-longtime}



\begin{quote}
Before the new one: say the Tamil for come again, then the Tamil for have you eaten.
\end{quote}

\subsection*{You'll want to know: \ta{ரொம்ப} \ta{நாளாச்சு}}
\textbf{\ta{ரொம்ப} \ta{நாளாச்சு}} (\emph{rompa nāḷāccu}) — \textquotedblleft{}it's been a long time\textquotedblright{}, in everyday speech.

\textbf{\ta{ரொம்ப}} is spoken \textquotedblleft{}very\textquotedblright{}; \textbf{\ta{நாள்}} is \textquotedblleft{}day\textquotedblright{}.

\begin{grammarlens}[title={What you've built}]
A greeting for someone you have not seen.
\end{grammarlens}

\subsection*{Your turn}
\begin{itemize}
\raggedright
  \item \emph{Say it:} \emph{rompa nāḷāccu}
  \item \emph{Say it:} \emph{rompa nāḷāccu}, once more
  \item \emph{Say it:} say \emph{rompa nāḷāccu!}
  \item \emph{Recall:} say the Tamil for come again, then the Tamil for have you eaten, then say \emph{rompa nāḷāccu} again
\end{itemize}

\begin{tcolorbox}[breakable,colback=teal!4,colframe=teal!35!black,title={Before you move on}]
What does \ta{ரொம்ப} \ta{நாளாச்சு} mean? (\textquotedblleft{}It's been a long time\textquotedblright{}.) Say it once more.
\end{tcolorbox}
\section[eṉṉa vicēṣam?]{\ta{என்ன} \ta{விசேஷம்}? (eṉṉa vicēṣam?) — what's new?}

\label{lesson:TA-C108-whatsnew}



\begin{quote}
Before the new one: say the Tamil for have you eaten, then the Tamil for it's been a long time.
\end{quote}

\subsection*{You'll want to know: \ta{என்ன} \ta{விசேஷம்}?}
\textbf{\ta{என்ன} \ta{விசேஷம்}?} (\emph{eṉṉa vicēṣam?}) — \textquotedblleft{}what's new?\textquotedblright{}, literally \textquotedblleft{}what's special?\textquotedblright{}.

\begin{grammarlens}[title={What you've built}]
A friendly opener.
\end{grammarlens}

\subsection*{Your turn}
\begin{itemize}
\raggedright
  \item \emph{Say it:} \emph{eṉṉa vicēṣam?}
  \item \emph{Say it:} \emph{eṉṉa vicēṣam?}, once more
  \item \emph{Say it:} ask \emph{eṉṉa vicēṣam?}
  \item \emph{Recall:} say the Tamil for have you eaten, then the Tamil for it's been a long time, then say \emph{eṉṉa vicēṣam?} again
\end{itemize}

\begin{tcolorbox}[breakable,colback=teal!4,colframe=teal!35!black,title={Before you move on}]
What does \ta{என்ன} \ta{விசேஷம்}? mean? (\textquotedblleft{}What's new?\textquotedblright{}.) Say it once more.
\end{tcolorbox}
\section[nallā irukkīṅkaḷā?]{\ta{நல்லா} \ta{இருக்கீங்களா}? (nallā irukkīṅkaḷā?) — are you well? (spoken)}

\label{lesson:TA-C108-areyouwell}



\begin{quote}
Before the new one: say the Tamil for it's been a long time, then the Tamil for what's new.
\end{quote}

\subsection*{You'll want to know: \ta{நல்லா} \ta{இருக்கீங்களா}?}
\textbf{\ta{நல்லா} \ta{இருக்கீங்களா}?} (\emph{nallā irukkīṅkaḷā?}) — \textquotedblleft{}are you well?\textquotedblright{}, as people say it.

\textbf{\ta{நல்லா}} is spoken \textbf{\ta{நன்றாக}}, \textquotedblleft{}well\textquotedblright{}.

\begin{grammarlens}[title={What you've built}]
The spoken shape of a question you know.
\end{grammarlens}

\subsection*{Your turn}
\begin{itemize}
\raggedright
  \item \emph{Say it:} \emph{nallā irukkīṅkaḷā?}
  \item \emph{Say it:} \emph{nallā irukkīṅkaḷā?}, once more
  \item \emph{Say it:} ask \emph{nallā irukkīṅkaḷā?}
  \item \emph{Recall:} say the Tamil for it's been a long time, then the Tamil for what's new, then say \emph{nallā irukkīṅkaḷā?} again
\end{itemize}

\begin{tcolorbox}[breakable,colback=teal!4,colframe=teal!35!black,title={Before you move on}]
What does \ta{நல்லா} \ta{இருக்கீங்களா}? mean? (\textquotedblleft{}Are you well? (spoken)\textquotedblright{}.) Say it once more.
\end{tcolorbox}
\section[eppaṭi irukkīṅka?]{\ta{எப்படி} \ta{இருக்கீங்க}? (eppaṭi irukkīṅka?) — how are you? (spoken)}

\label{lesson:TA-C108-howareyou}



\begin{quote}
Before the new one: say the Tamil for what's new, then the Tamil for are you well.
\end{quote}

\subsection*{You'll want to know: \ta{எப்படி} \ta{இருக்கீங்க}?}
\textbf{\ta{எப்படி} \ta{இருக்கீங்க}?} (\emph{eppaṭi irukkīṅka?}) — \textquotedblleft{}how are you?\textquotedblright{}, the everyday spoken form.

It is the written \textbf{\ta{எப்படி} \ta{இருக்கிறீர்கள்}?} worn short in speech.

\begin{grammarlens}[title={What you've built}]
Five greetings people actually use.
\end{grammarlens}

\subsection*{Your turn}
\begin{itemize}
\raggedright
  \item \emph{Say it:} \emph{eppaṭi irukkīṅka?}
  \item \emph{Say it:} \emph{eppaṭi irukkīṅka?}, once more
  \item \emph{Say it:} ask \emph{eppaṭi irukkīṅka?}
  \item \emph{Recall:} say the Tamil for what's new, then the Tamil for are you well, then say \emph{eppaṭi irukkīṅka?} again
\end{itemize}

\begin{tcolorbox}[breakable,colback=teal!4,colframe=teal!35!black,title={Before you move on}]
What does \ta{எப்படி} \ta{இருக்கீங்க}? mean? (\textquotedblleft{}How are you? (spoken)\textquotedblright{}.) Say it once more.
\end{tcolorbox}
\section[(dialogue)]{Hello and goodbye, the everyday way — a second pass}

\label{lesson:TA-R108-second-pass-leave-greet}



\begin{quote}
You meet an old friend around lunchtime. Say two things.
\end{quote}

\begin{grammarlens}[title={Grammar Lens: written and spoken}]
\noindent
\begin{tabularx}{\linewidth}{@{}>{\raggedright\arraybackslash}X>{\raggedright\arraybackslash}X@{}}
\toprule
\textbf{spoken} & \textbf{written} \\
\midrule
\textbf{\ta{நல்லா}} & \textbf{\ta{நன்றாக}} \\
\textbf{\ta{இருக்கீங்க}} & \textbf{\ta{இருக்கிறீர்கள்}} \\
\bottomrule
\end{tabularx}

Spoken Tamil wears the written forms short. Both are right; they are for different places.
\end{grammarlens}

\subsection*{Your turn}
Say these aloud:

\begin{itemize}
\raggedright
  \item ask a friend how they are, the spoken way; then what's new
  \item at the door: shall we leave? see you later
  \item wish a traveller a good journey; ask them to come again
  \item take your leave formally
\end{itemize}

\begin{tcolorbox}[breakable,colback=teal!4,colframe=teal!35!black,title={Before you move on}]
I take my leave? (\textbf{\ta{விடைபெறுகிறேன்}}.) Shall we leave? (\textbf{\ta{கிளம்பலாமா}?}) See you later? (\textbf{\ta{அப்புறம்} \ta{பார்க்கலாம்}}.) A good journey? (\textbf{\ta{இனிய} \ta{பயணம்}}.) Come again? (\textbf{\ta{மீண்டும்} \ta{வாருங்கள்}}.) Have you eaten? (\textbf{\ta{சாப்பிட்டீர்களா}?}) Long time? (\textbf{\ta{ரொம்ப} \ta{நாளாச்சு}}.) What's new? (\textbf{\ta{என்ன} \ta{விசேஷம்}?}) Are you well? (\textbf{\ta{நல்லா} \ta{இருக்கீங்களா}?}) How are you? (\textbf{\ta{எப்படி} \ta{இருக்கீங்க}?})
\end{tcolorbox}
\section[(dialogue)]{Leaving well — a second pass}

\label{lesson:TA-R108-second-pass-care}



\begin{quote}
It has got late. Say so, and suggest you all go.
\end{quote}

\subsection*{Your turn}
Say these aloud:

\begin{itemize}
\raggedright
  \item tell your host to take care; they tell you to go safely
  \item say tomorrow is a holiday
  \item next week, greet the same host: are you well? how are you?
\end{itemize}

\begin{tcolorbox}[breakable,colback=teal!4,colframe=teal!35!black,title={Before you move on}]
Take care? (\textbf{\ta{பார்த்துக்கொள்ளுங்கள்}}.) Safely? (\textbf{\ta{பத்திரமாக}}.) It's late? (\textbf{\ta{நேரமாகிவிட்டது}}.) Let's go? (\textbf{\ta{போகலாம்}}.) A holiday? (\textbf{\ta{விடுமுறை}}.) Are you well? (\textbf{\ta{நல்லா} \ta{இருக்கீங்களா}?}) How are you? (\textbf{\ta{எப்படி} \ta{இருக்கீங்க}?})
\end{tcolorbox}
